% Part: computability
% Chapter: recursive-functions
% Section: primitive-recursion

\documentclass[../../../include/open-logic-section]{subfiles}

\begin{document}

\olfileid{cmp}{rec}{pre}
\olsection{بازگشت اولیه}

یکی از ویژگی‌های اعداد طبیعی این است که با به‌کاربردن متناهیِ عمل
جانشین~$+1$، می‌توان از~$0$ به هر عدد طبیعی رسید---هر عدد طبیعی یا
خود~$0$ است یا جانشینِ \dots{} جانشینِ~$0$. یک راه تعیین تابعی
$h\colon\Nat \to \Nat$ با بهره‌گیری از این واقعیت چنین است:
(الف)~مقدار $h$ را برای آرگومان~$0$ تعیین کنیم و (ب)~نیز مشخص کنیم
که با داشتن مقدار $h(x)$ چگونه مقدار $h(x+1)$ را محاسبه کنیم. بند
(الف) مستقیماً می‌گوید $h(0)$ چیست، پس $h$ برای~$0$ تعریف شده است.
اکنون با دستورِ بند (ب) برای $x=0$ می‌توانیم $h(1)=h(0+1)$ را از
$h(0)$ محاسبه کنیم. با همان دستور برای $x=1$، مقدار
$h(2)=h(1+1)$ را از~$h(1)$ محاسبه می‌کنیم و به همین ترتیب ادامه
می‌دهیم. برای هر عدد طبیعی~$x$ سرانجام به گامی می‌رسیم که در آن
$h(x+1)$ را از~$h(x)$ تعریف می‌کنیم؛ ازاین‌رو $h(x)$ برای همهٔ
$x\in\Nat$ تعریف شده است.

برای نمونه، فرض کنید $h\colon \Nat \to \Nat$ را با دو معادلهٔ زیر
تعیین کنیم:
\begin{align*}
h(0) & =  1\\
h(x+1) & =  2 \cdot h(x)
\end{align*}
اگر از پیش ضرب‌کردن را بدانیم، این معادله‌ها اطلاعات لازم برای بندهای
(الف) و~(ب) بالا را می‌دهند. با کاربرد پیاپی معادلهٔ دوم به دست
می‌آوریم:
\begin{align*}
  h(1) & = 2\cdot h(0) = 2,\\
  h(2) & = 2\cdot h(1) = 2\cdot 2,\\
  h(3) & = 2 \cdot h(2) = 2\cdot 2 \cdot 2,\\
  & \vdots
\end{align*}
می‌بینیم تابع~$h$ که تعیین کرده‌ایم همان $h(x)=2^x$ است.

ویژگی یادشدهٔ اعداد طبیعی تضمین می‌کند که تنها یک تابع~$h$ این دو
شرط را برآورده می‌کند. جفتی از معادله‌ها مانند این‌ها
\emph{تعریف تابع~$h$ با بازگشت اولیه} نام دارد. این تعریف
«بازگشتی» است، زیرا خودِ $h$ در تعریف---به‌ویژه در سمت راست معادلهٔ
دوم---حضور دارد. «اولیه» است، زیرا برای تعریف $h(x+1)$ تنها از مقدار
$h(x)$، یعنی مقدار بلافاصله پیشین، استفاده می‌کنیم. این ساده‌ترین راه
تعریف بازگشتی تابعی روی~$\Nat$ است.

حتی توابع بنیادی‌تری مانند جمع و ضرب را می‌توان با بازگشت اولیه تعریف
کرد. بااین‌حال، در این موارد توابع مورد نظر دوموضعی‌اند. یکی از مواضع
آرگومان را ثابت می‌گیریم و از دیگری برای بازگشت استفاده می‌کنیم. برای
مثال، برای تعریف $\Add(x,y)$ می‌توانیم~$x$ را ثابت بگیریم، نخست مقدار
را برای $y=0$ تعریف کنیم و سپس مقدار مربوط به $y+1$ را برحسب~$y$
به دست دهیم. چون $x$ ثابت است، در هر دو سوی معادله‌های تعریف ظاهر
می‌شود:
\begin{align*}
\Add(x,0) & =  x\\
\Add(x,y+1) & =  \Add(x,y)+1
\end{align*}
این معادله‌ها مقدار $\Add$ را برای همهٔ $x$ها \emph{و} همهٔ $y$ها
تعیین می‌کنند. برای یافتن $\Add(2,3)$، مثلاً، معادله‌های تعریف را با
$x=2$ به کار می‌بریم: از نخستین معادله $\Add(2,0)=2$ را می‌یابیم و
سپس با معادلهٔ دوم به‌ترتیب
$\Add(2,1)=2+1=3$، $\Add(2,2)=3+1=4$ و $\Add(2,3)=4+1=5$ را به
دست می‌آوریم.

در تعریف $\Add$ در سمت راست معادلهٔ دوم از~$+$ استفاده کردیم، اما
تنها برای افزودن~$1$. به بیان دیگر، تابع جانشین
$\Succ(z)=z+1$ را به کار بردیم و آن را بر مقدار پیشین
$\Add(x,y)$ اعمال کردیم تا $\Add(x,y+1)$ را تعریف کنیم. پس می‌توان
تعریف بازگشتی را برحسب تابعی یگانه دید که بر مقدار پیشین اعمال
می‌شود. بااین‌همه، زیانی ندارد---و گاه ضروری است---که بگذاریم این تابع
نه‌فقط به مقدار پیشین، بلکه به $x$ و~$y$ نیز وابسته باشد. در نظر
بگیرید:
\begin{align*}
  \Mult(x,0) & =  0 \\
  \Mult(x,y+1) & =  \Add(\Mult(x,y),x)
\end{align*}
این تعریفی با بازگشت اولیه برای تابع $\Mult$ است که تابع $\Add$ را
هم بر مقدار پیشین $\Mult(x,y)$ و هم بر آرگومان نخست~$x$ اعمال
می‌کند. این معادله‌ها نیز تابع~$\Mult(x,y)$ را برای همهٔ آرگومان‌های
$x$ و~$y$ تعریف می‌کنند. برای نمونه، $\Mult(2,3)$ با محاسبهٔ پیاپی
$\Mult(2,0)$، $\Mult(2,1)$، $\Mult(2,2)$ و~$\Mult(2,3)$ تعیین
می‌شود:
\begin{align*}
  \Mult(2,0) & = 0\\
  \Mult(2,1) & = \Mult(2,0+1) =
  \Add(\Mult(2,0), 2) = \Add(0, 2) = 2\\
  \Mult(2,2) & = \Mult(2,1+1) =
  \Add(\Mult(2,1), 2) = \Add(2, 2) = 4\\
  \Mult(2,3) & = \Mult(2,2+1) =
  \Add(\Mult(2,2), 2) = \Add(4, 2) = 6
\end{align*}

پس الگوی کلی چنین است: برای تعریف تابع
$h(x_0,\dots,x_{k-1},y)$ با بازگشت اولیه، دو معادله به دست می‌دهیم.
معادلهٔ نخست مقدار $h(x_0,\dots,x_{k-1},0)$ را بدون ارجاع به~$h$
تعریف می‌کند. معادلهٔ دوم مقدار $h(x_0,\dots,x_{k-1},y+1)$ را
برحسب $h(x_0,\dots,x_{k-1},y)$، آرگومان‌های دیگر
$x_0$، \dots،~$x_{k-1}$ و~$y$ تعریف می‌کند. در معادلهٔ دوم تنها
می‌توان از مقدار بلافاصله پیشینِ~$h$ استفاده کرد. اگر عمل‌های سمت
راست این دو معادله را خود تابع‌هایی چون $f$ و~$g$ در نظر بگیریم،
الگوی کلی تعریف تابع تازهٔ~$h$ با بازگشت اولیه چنین است:
\begin{align*}
  h(x_0, \dots, x_{k-1}, 0) & = f(x_0, \dots, x_{k-1})\\
  h(x_0, \dots, x_{k-1}, y+1) & =
  g(x_0, \dots, x_{k-1}, y, h(x_0, \dots, x_{k-1}, y))
\end{align*}
در مورد $\Add$، داریم $k=1$ و $f(x_0)=x_0$ (تابع همانی)، و
$g(x_0,y,z)=z+1$ (تابع سه‌موضعی‌ای که جانشین آرگومان سومش را
بازمی‌گرداند):
\begin{align*}
  \Add(x_0, 0) & = f(x_0) = x_0\\
  \Add(x_0, y+1) & = g(x_0, y, \Add(x_0, y)) =
  \Succ(\Add(x_0, y))
\end{align*}
در مورد $\Mult$، داریم $f(x_0)=0$ (تابع ثابتی که همواره~$0$ را
بازمی‌گرداند) و $g(x_0,y,z)=\Add(z,x_0)$ (تابع سه‌موضعی‌ای که مجموع
آرگومان آخر و نخست خود را بازمی‌گرداند):
\begin{align*}
  \Mult(x_0,0) & =  f(x_0) = 0 \\
  \Mult(x_0,y+1) & =  g(x_0, y, \Mult(x_0,y)) =
  \Add(\Mult(x_0,y), x_0)
\end{align*}

\end{document}
